\documentclass[11pt]{article}

\usepackage[a4paper,margin=1in]{geometry}
\usepackage{amsmath,amssymb,amsthm}
\usepackage[T1]{fontenc}
\usepackage{lmodern}
\usepackage{hyperref}

\newtheorem{theorem}{Theorem}
\newtheorem{lemma}[theorem]{Lemma}
\newcommand{\Z}{\mathbb Z}

\title{The Exact Value of the Zero-Sum Constant $s_1(4)$}
\author{Xiaojian Zeng\\
PBC School of Finance, Tsinghua University, Beijing 100190, China\\
\texttt{xiaojianzeng@foxmail.com}}
\date{}

\begin{document}

\maketitle

\begin{abstract}
For a positive integer $k$, consider sequences $a_1,\ldots,a_k$ of
integers not divisible by $4$.  The constant $s_1(4)$ is the least $k$
for which every such sequence has four distinct positions whose sum is
divisible by $4$ but not by $16$.  We prove that $s_1(4)=9$.  The lower
bound is witnessed by $(1,1,1,2,2,2,10,13)$.  For the upper bound, the
problem is reduced without loss to the $167{,}960$ multiplicity vectors
of size nine on the twelve allowed residue classes modulo $16$.  A
bounded-multiplicity generating-function dynamic program computes the
attainable four-term residues using exact Boolean arithmetic and finds no
avoiding vector.  Two structurally different exact implementations give
the same result.  This determines only the case $n=4$ and does not address
the general conjecture for $s_1(n)$.
\end{abstract}

\section{Statement and context}

Sun introduced $s_r(n)$ by requiring a subsequence of exactly $n$
positions, rather than a nonempty subsequence or one having merely a
bounded length \cite{Sun2026ZeroSum}.  Specializing the definition to
$r=1$ and $n=4$, let $P(k)$ be the assertion that every ordered sequence
$a_1,\ldots,a_k$ with $4\nmid a_i$ has distinct indices
$i_1<i_2<i_3<i_4$ such that
\[
  4\mid a_{i_1}+a_{i_2}+a_{i_3}+a_{i_4}
  \quad\text{and}\quad
  16\nmid a_{i_1}+a_{i_2}+a_{i_3}+a_{i_4}.
\]
Then $s_1(4)=\min\{k\geq 1:P(k)\}$.  Equal integer values at different
positions may all be selected; it is the positions that must be distinct.

Sun proved $2nr+1\leq s_r(n)$ for $n>2$ and, using the equality
$s(\Z_n^2)=4n-3$ established through the resolution of Kemnitz's
conjecture \cite{Reiher2007Kemnitz}, obtained $s_1(n)\leq 4n-3$ for
$n\geq4$ \cite{Sun2026ZeroSum}.  Thus the previously supplied bounds in
the present case were
\[
  9\leq s_1(4)\leq13,
\]
and Sun conjectured $s_1(n)=2n+1$ for every $n>2$.  Gao, Jiang, and Mu
subsequently determined related local constants in which one may choose
some nonempty number of terms or at most $n$ terms
\cite{GaoJiangMu2026Local}; those quantifiers do not give the
exactly-four assertion considered here.

The contribution below is a complete determination of the frozen
$n=4$ case.  The lower bound is proved directly, independently of the
general lower bound just cited, and the upper bound is an exact finite
proof.

\section{Exact reduction modulo 16}

Put
\[
 A=\{1,2,3,5,6,7,9,10,11,13,14,15\}\subset\Z/16\Z,
 \qquad T=\{4,8,12\}.
\]

\begin{lemma}[Residue and multiplicity reduction]\label{lem:reduction}
The assertion $P(k)$ holds if and only if every multiplicity vector
$(c_r)_{r\in A}\in\Z_{\geq0}^{12}$ with $\sum_{r\in A}c_r=k$ contains
nonnegative integers $q_r\leq c_r$ such that
\[
  \sum_{r\in A}q_r=4
  \quad\text{and}\quad
  \sum_{r\in A}q_r r\pmod {16}\in T.
\]
There are $\binom{k+11}{11}$ such multiplicity vectors.
\end{lemma}

\begin{proof}
An integer not divisible by $4$ has residue in $A$, and a residue modulo
$16$ is divisible by $4$ but nonzero modulo $16$ precisely when it lies
in $T$.  Thus both the hypothesis and the desired conclusion depend only
on the entries modulo $16$; this also covers negative integers.

Permuting the sequence positions does not affect whether a suitable set
of four positions exists.  Each permutation orbit is represented by the
unique vector of multiplicities of the twelve residues in $A$, and every
such vector expands to an ordered sequence.  Selecting positions is
equivalent to choosing bounded submultiplicities $0\leq q_r\leq c_r$.
Finally, stars and bars gives $\binom{k+11}{11}$ vectors.
\end{proof}

At lengths eight and nine the respective domain sizes are
\[
 \binom{19}{11}=75{,}582,
 \qquad
 \binom{20}{11}=167{,}960.
\]

\section{The sharp lower bound}

\begin{lemma}\label{lem:lower}
The sequence
\[
  (1,1,1,2,2,2,10,13)
\]
has no four positions whose sum modulo $16$ belongs to $T$.
Consequently $s_1(4)\geq9$.
\end{lemma}

\begin{proof}
Suppose a four-position selection uses $a,b,c,d$ copies of
$1,2,10,13$, respectively.  Then
\[
 a+b+c+d=4,\qquad a,b\leq3,\qquad c,d\leq1,
\]
and, after eliminating $a$, its sum is
\[
 a+2b+10c+13d\equiv4+b+9c+12d\pmod {16}.
\]
For $(c,d)=(0,0),(1,0),(0,1),(1,1)$, the complete attainable residue
sets are, respectively,
\[
 \{5,6,7\},\quad
 \{13,14,15,0\},\quad
 \{0,1,2,3\},\quad
 \{9,10,11\}.
\]
Their union omits $4,8,12$.  Hence $P(8)$ fails.  Since $P(k)$ implies
$P(k+1)$ by applying $P(k)$ to any $k$ positions, the least valid length
cannot be at most eight.
\end{proof}

\section{Exhaustive upper bound}

For a fixed multiplicity vector $c$, order the elements of $A$ as
$r_1,\ldots,r_{12}$.  Let $R_{i,j}$ be the set of residues modulo $16$
obtainable by choosing exactly $j$ positions from the first $i$ residue
classes, respecting their multiplicities.  Initialize
\[
 R_{0,0}=\{0\},\qquad R_{0,j}=\varnothing\quad(1\leq j\leq4),
\]
and update by
\begin{equation}\label{eq:dp}
 R_{i+1,j}=
 \bigcup_{0\leq q\leq\min(c_{r_{i+1}},j)}
 \bigl(R_{i,j-q}+q r_{i+1}\bigr)\pmod {16}.
\end{equation}
Only degrees $0,1,2,3,4$ are retained.  Each set is represented by a
$16$-bit mask, so every operation and comparison is exact.

\begin{lemma}[Exhaustive certificate]\label{lem:certificate}
For every multiplicity vector $c\in\Z_{\geq0}^{12}$ with
$\sum_{r\in A}c_r=9$, the final set $R_{12,4}$ intersects $T$.
\end{lemma}

\begin{proof}
Induction on $i$ proves the invariant defining $R_{i,j}$.  Indeed, if
exactly $q$ copies of residue $r_{i+1}$ are selected, the remaining
$j-q$ terms contribute a member of $R_{i,j-q}$, and every bounded choice
appears in exactly one of the corresponding cases of
\eqref{eq:dp}.  Discarding degrees greater than four cannot affect degree
four because selected cardinalities never decrease.

It remains to cover the finite domain.  The primary verifier recursively
assigns each of the first eleven coordinates every value from zero to the
current remainder and forces the twelfth coordinate to equal the final
remainder.  Each weak composition of nine has one and only one such
root-to-leaf path.  The recursion visited $167{,}960=\binom{20}{11}$
leaves.  Applying \eqref{eq:dp} at every leaf found zero vectors for which
$R_{12,4}\cap T=\varnothing$.

Two different implementations checked the finite conclusion.  A direct
verifier enumerates nondecreasing residue multisets and tests their
four-position subsets.  A third verifier generates multiplicity vectors
through the separator form of stars and bars, independently generates
the $274$ admissible four-term multiplicity patterns, and tests componentwise
domination.  The following results were obtained with exact integer and
Boolean arithmetic:
\begin{center}
\begin{tabular}{lrr|rr}
\hline
& \multicolumn{2}{c|}{length $8$} & \multicolumn{2}{c}{length $9$}\\
verifier & tested & avoiders & tested & avoiders\\
\hline
multiplicity dynamic program & $75{,}582$ & $160$ & $167{,}960$ & $0$\\
direct four-position enumeration & $75{,}582$ & $160$ & $167{,}960$ & $0$\\
pattern-domination verifier & $75{,}582$ & $160$ & $167{,}960$ & $0$\\
\hline
\end{tabular}
\end{center}
The three routes use neither floating-point arithmetic nor randomized
decisions.  The two frozen verifier sources and their recorded outputs
accompany the paper; the third reconstruction is documented in the
mathematical audit.  This establishes the claimed universal finite
statement.
\end{proof}

\begin{theorem}
Under the definition above,
\[
  s_1(4)=9.
\]
\end{theorem}

\begin{proof}
Lemma~\ref{lem:lower} gives $s_1(4)\geq9$.  By
Lemma~\ref{lem:reduction}, Lemma~\ref{lem:certificate} says that every
allowed sequence of nine integers has the required four positions, so
$s_1(4)\leq9$.
\end{proof}

\section{Reproducibility, scope, and limitations}

The primary program implements weak-composition enumeration and the
bounded-multiplicity recurrence \eqref{eq:dp}.  The second program instead
uses sorted residue multisets and explicit four-position combinations;
the third uses a separator bijection and pattern domination.  Each source
reconstructs its own finite objects rather than accepting an externally
generated list.  The primary run used Python 3.10.12 and only the standard
library.  The decisive claims are the domain bijection, the recurrence
invariant, and the exact exhaustive outputs, not a numerical fit or a
successful heuristic search.

This certificate architecture is comparable to the independently
reconstructed integer certificates used for other finite combinatorial
models in \cite{ZengEtAl2026FiniteCertificates}; that work is cited only
for the reproducibility methodology and is not a premise of the theorem
above.

The computation proves a theorem about arbitrary integer sequences
because Lemma~\ref{lem:reduction} is a lossless finite reduction.  By
contrast, it supplies no conclusion for $s_1(n)$ when $n\neq4$, for ranks
$r\neq1$, or for variants permitting fewer or more than four selected
positions.  A bounded literature search through 10 September 2026 found
Sun's conjecture and the distinct local constants of Gao, Jiang, and Mu,
but no equivalent determination of this exact case.  This search record
is not a claim of priority.

\paragraph{Computational and AI assistance.}
OpenAI Codex assisted with literature searching, the design and
implementation of the exhaustive checks, adversarial audit organization,
and initial manuscript preparation.  Language-model output is not used as
mathematical evidence.  The finite conclusion rests on the explicit
reduction, exact deterministic programs, recorded outputs, and a separate
reconstruction of the decisive computation.  No proof assistant was used.

\bibliographystyle{plain}
\bibliography{references}

\end{document}
